\cotefile{157-4}
\scanpage{348}
% @ head of sheet
\mspage{72}
% @ top, lines 1–4
une \inserted{(petite)} catégorie \unsure{filtrante} quelconque, et quitte à la \unsure{remplacer} par une catégorie \unsure{cofinale}, \unsure{pps|sps|p. ex.} $I$ un ensemble \unsure{ordonné}. L'axiome $\mathrm{Loc}(8\ \mathrm{bis})$ correspond au cas où $I=\mathbb{N}$.

% @ upper middle, lines 5–8
Notons que grâce à la construction de \unsure{la|l'é} \deleted{\unsure{équi|qu}} catégorie des chemins, nous avons une \deleted{\unsure{construction}} \inserted{factorisation}, fonctorielle en $f$, de \deleted{\unsure{toutes}} flèches $f\inserted{{:}\,X\to Y}$ de $\mathrm{Cat}$ en $f=pi$ :
% @ centre, first factorization
\diagram{\begin{tikzcd}[ampersand replacement=\&]
X \arrow[r,"i=i_f"] \& \overline X(f)=\overline X \arrow[r,"p=p_f"] \& Y{,}
\end{tikzcd}}
% @ middle, lines 10–12
avec $i\in W_{\unsure{\omega|W|w}}$, $p\in\mathrm{Parf}_{W_\omega}$, a fortiori $i\in W$, $p\in\mathrm{Parf}_W$. Ce foncteur $f\longrightarrow(\overline X(f),p_f)$ \unsure{commute} aux limites inductives.
% @ lower middle, lines 13–16
Ainsi, si on applique ceci à \unsure{chacun} des $f_\alpha:X_\alpha\longrightarrow Y_\alpha$, les factorisant en
\diagram{\begin{tikzcd}[ampersand replacement=\&]
X_\alpha \arrow[r,"i_\alpha",dash] \& \overline X_\alpha \arrow[r,"p_\alpha",dash] \& Y_\alpha{,}
\end{tikzcd}}
% @ foot, lines 17–21
\deleted{\unsure{on}} \unsure{on} trouve un système inductif de diagrammes de type $\Delta^2$ (factorisant les $f_i$), et en passant à la $\underrightarrow{\lim}$ on trouve
\diagram{\begin{tikzcd}[ampersand replacement=\&]
X \arrow[r,"i_{\deleted{\unsure{\alpha|\omega}}}"] \& \overline X \arrow[r,"p",dash] \& Y{,}
\end{tikzcd}}
qui n'est autre que la factorisation \unsure{canonique}
% archive sheet number: 347
\scanpage{349}
% @ head of sheet
\mspage{73}
% @ top, lines 1–5
de $f=\underrightarrow{\lim} f_i$ en $pi$, avec $i\in W_{\unsure{\omega|W|w}}$, $p\in\mathrm{Parf}_{W_\omega}$. (NB Alors que ni $W_\omega$, ni $\mathrm{Parf}_{W_\omega}$ ne sont stables par $\underrightarrow{\lim}$ filtrantes, à coup sûr !). On a donc $i\in W$, et \deleted{\unsure{pour}} la condition $f\in W$ i.e. $pi\in W$, équivaut donc à $p\in W$.
% @ upper middle, lines 6–8
Cela nous ramène, pour vérifier l'axiome $\mathrm{Loc}\,8_{\unsure{\alpha|a}}$ (resp $\mathrm{Loc}\,8\mathrm{bis}$), au cas où les $f_i$ sont dans $\mathrm{Parf}_{W_\omega}$.

% @ centre, lines 9–13
Supposons maintenant que $W$ satisfasse à $D_{W_\omega}$ (la plus faible de toutes les variantes \inserted{envisagées ici} de l'axiome des $W$-fibrations \inserted{$\mathrm{Loc}(7)$}, à la seule exception de $F$), \deleted{Alors} \deleted{et de plus à} Alors on a
% @ below centre, lines 14–15
\[
p\in W \quad\Longleftrightarrow\quad
\begin{aligned}&\text{les fibres de }p\text{ sont}\\[-2pt]
& W\text{-asphériques}\end{aligned},
\]
% @ lower third, lines 16–22
l'implication $\Longrightarrow$ résultant de l'hypothèse $D_{W_\omega}$ $\mathrm{Parf}_{W_\omega}\cap W\subset W^{\unsure{u|n}}$, et $\Longleftarrow$ résultant du \unsure{critère|crit.} des morphismes $W$-lisses etc, via $\mathrm{Loc}(3)$. De $\unsure{=|\text{même}}$, l'hypothèse que les $f_i\in W$ équivaut à celle que les fibres de $p_i$ \unsure{sont} $W$-asphériques. \deleted{\unsure{Ceci}} D'autre part, il est immédiat que
% archive sheet number: 348
\scanpage{350}
% @ head of sheet
\mspage{74}
% @ top, lines 1–2
toute fibre de $p$ est $\inserted{\unsure{\sim|\simeq}}\underrightarrow{\lim}_{I}$ des fibres correspondantes des $p_\alpha$. De cette façon, on trouve

% @ upper third, lines 3–6
\uline{Proposition} 5.1. Supposons que $W$ satisfasse l'axiome \inserted{des $W$-fibrations} $\mathrm{Loc}(7\,D_{W_{\unsure{\omega|W|w}}})$ (ce qui est donc le cas s'il satisfait un des axiomes \inserted{plus forts} $\mathrm{Loc}(7\,a\ \text{ou}\ b\ \text{ou}\ c)$).

% @ middle, lines 7–12
Pour que $W$ satisfasse à $\mathrm{Loc}(8)$ (resp. à $\mathrm{Loc}(8\mathrm{bis})$), il f. et s. \deleted{qu'il} \deleted{soit} que l'ensemble des objets de $\mathrm{Cat}$ qui sont $W$-asphériques soit stable par $\underrightarrow{\lim}$ filtrantes (resp. par $\underrightarrow{\lim}$ \deleted{\unsure{f} \unsure{filtrantes}} de suites \unsure{dénombrables}).

% @ lower third, lines 13–16
\deleted{\unsure{Corollaire}} \inserted{5.2.} Notons que la catégorie $\mathrm{Cat}$ est un \unsure{paratopos|prétopos}, - a fortiori elle est \uline{accessible}, \inserted{cf $\overline{\underline{\mathrm{XVIII}}}$ \S\,8),} et il en est donc de $=$ de $\underline{\mathrm{Fl}}(\mathrm{Cat})$. Ainsi on a dans $\mathrm{Fl}(\mathrm{Cat})=\mathrm{Ob}\,\underline{\deleted{\unsure{\mathrm{E}}}\mathrm{Fl}}(\mathrm{Cat})$ la notion
% @ foot, lines 17–21
de \uline{partie accessible}. On dira donc qu'une partie $W$ de $\mathrm{Fl}\,\mathrm{Cat}$ (plus gén. d'une catégorie $\mathrm{Fl}(M)$, où $M$ est une catégorie accessible) est \uline{accessible}, si elle l'est \deleted{dans la} \unsure{p.r.} à la catégorie ambiante
% archive sheet number: 349
\scanpage{351}
% @ head of sheet
\mspage{75}
% @ top, lines 1–2
$\underline{\mathrm{Fl}}\,\mathrm{Cat}$. Ceci posé, l'axiome $\mathrm{Loc}\,9$ est le suivant

% @ upper third, lines 3–4
\[
\begin{aligned}
&\underline{\mathrm{Loc}}\ 9:\quad W\text{ est une partie \uline{accessible} de }\mathrm{Fl}(\mathrm{Cat})\\
&=\mathrm{Ob}\,\underline{\mathrm{Fl}}(\mathrm{Cat}).
\end{aligned}
\]
% @ middle, lines 5–11
5.3. Rappelons que cette notion un peu technique se simplifie considérablement, si on suppose que $W$ satisfait l'axiome des limites $\mathrm{Loc}(8)$. Alors $W$ est accessible \unsure{sss|ssi} il existe une \uline{petite} partie $W_{\unsure{\bullet|0|\circ}}$ de $\mathrm{Fl}(\mathrm{Cat})$, telle que $W$ soit la ``$L_1$-enveloppe'' de $W_0$, i.e.
% @ lower middle, lines 12–16
la plus petite partie de $\mathrm{Fl}(\mathrm{Cat})$ contenant $W_\bullet$ et \inserted{$L_1$-}stable \deleted{par} (i.e stable par $\underrightarrow{\lim}$ filtrantes). \deleted{\illegible{1}} S'il en est ainsi, \inserted{il est clair que} on peut \deleted{\unsure{trouver un toujours}} \inserted{prendre $W_\bullet$ de} la forme
% @ lower third, cardinal formula and right-hand insertion
\[
W_{\unsure{\pi|\kappa}}=W\cap\mathrm{Fl}_{\pi}(\mathrm{Cat})
\quad \inserted{\begin{aligned}&\text{où }\unsure{\kappa|\pi}\text{ est un}\\&\text{cardinal infini}\\&\text{convenable}\end{aligned}}
\]
% @ foot, last four lines
où $\mathrm{Fl}_{\kappa}(\mathrm{Cat})\overset{\mathrm{def}}{=}\mathrm{Fl}(\underline{\mathrm{Cat}}_{\pi})$, $\underline{\mathrm{Cat}}_{\pi}$ désignant la sous-catégorie \inserted{\unsure{str.}} pleine de $\mathrm{Cat}$ formée des \deleted{catégories} objets $X$ tels que $\operatorname{card} X\leq\pi$ (et par définition, $\operatorname{card} X=\operatorname{card}\mathrm{Fl}\,X$).
% archive sheet number: 350
\scanpage{352}
% @ head of sheet
\mspage{76}
% @ upper third, lines 1–6
Notons que la filtration qu'on vient de décrire sur $\mathrm{Cat}$ (indexée par les \unsure{petits} cardinaux infinis) n'est autre que sa filtration canonique, \deleted{\unsure{pour}} i.e. $\underline{\mathrm{Cat}}_{\unsure{\pi|\kappa}}$ est la sous-catégorie pleine de $\mathrm{Cat}$ formée des objets $\unsure{\kappa|\pi}$-accessibles \inserted{de $\mathrm{Cat}$}, et de $=$ dans $\underline{\mathrm{Fl}}(\mathrm{Cat})$. On
% @ upper middle, lines 7–9
\unsure{trouve} ainsi, grâce aux développements de XVIII \S\,13, la caractérisation suivante, que je rappelle ici pour mémoire \inserted{(en me bornant pour simplifier au cas où $W$ satisfasse à $\mathrm{Loc}(8)$)} :

% @ centre, beside the proposition
% slanted
\margin{\unsure{13.35|13.3.5} \unsure{cf|de} loc. cit}
% @ centre, proposition and cancelled first formulation
\uline{Prop.} 5.3.1. \deleted{Pour que $W$ soit accessible, il faut} \inserted{Soit $W$ une partie de $\mathrm{Fl}(\mathrm{Cat})$ stable \unsure{p'r|par} $\underrightarrow{\lim}$} filtrantes, et \unsure{soit|suffit}, pour \deleted{qu'il existe un cardinal \inserted{infini} $\pi$ ayant la propriété suivante : $W$ \deleted{satisfait} est \inserted{$L_\pi$-}stable \unsure{dans} $\mathrm{Fl}\,\mathrm{Cat}$ (i.e. stable par $\underrightarrow{\lim}$ grandes devant $\pi$) et de plus, pour \deleted{\unsure{tt}} système inductif $(f_i)_{i\in I}$ grd devant $\pi$} tout cardinal \inserted{infini} $\pi$,
% @ lower middle, definitions after the cancellation
$W_\pi=W\cap\mathrm{Fl}_{\pi}(\mathrm{Cat})$, et \inserted{$\underline W$,} $\underline W_\pi$ les sous-catégories \deleted{\unsure{pl.}} pleines correspondantes de $\mathrm{Fl}(\mathrm{Cat})$. On dit que $W$ est \uline{$\pi$-accessible} si $W$ est la $L$-enveloppe de $W_\pi$ (donc $W$ est accessible \unsure{ssi|est|cst} $\exists$ cardinal infini $\pi$ tel que $W$ soit $\pi$-accessible). Ceci posé, les conditions suivantes sont \unsure{équivalentes}
% @ foot, conditions a and b
a) $W$ est \deleted{\unsure{L}} $\pi$-accessible

b) Pour \unsure{tt} système inductif $(f_i)_{i\in I}$
% archive sheet number: 351
\scanpage{353}
% @ head of sheet
\mspage{77}
% @ top, continuation of b, lines 1–4
$L_1$-$\unsure{\kappa|\pi}$-adapté dans $\underline{\mathrm{Fl}}\,\mathrm{Cat}$ (cf ci-dessous), de limite inductive $f$, on a $f\in W$ \unsure{sss|ssi} la sous-\inserted{catégorie pl.} $I_W$ des $i\in\mathrm{Ob}\,I$ tels que $f_i\in \unsure{w|W}$ est cofinale dans $I$.

% @ upper third, lines 5–8
b$'$) (Quand on a choisi, pour \unsure{tte} $\deleted{\unsure{x}}f\in\deleted{\unsure{O}}\mathrm{Fl}(\mathrm{Cat})$, une représentation $L_1$-$\kappa$-adaptée de $f$ comme $\underrightarrow{\lim}$) : Si $f\in W$, alors la \deleted{\unsure{cat.}}\inserted{sous-cat. pl.} $I_W$ \unsure{définie} dans b) est cofinale dans $I$.

% @ centre, lines 9–10
Rappelons ici la notion de syst. inductif $L_1$-$\kappa$-adapté. Il s'agit des conditions :

% @ centre, conditions 1–3
1\textsuperscript{o}) $I$ est \unsure{grde|grande} devant $\kappa$, et stable par $\underrightarrow{\lim}$ filtrantes de cardinal $\leq\kappa$.

2\textsuperscript{o}) Le foncteur $i\mapsto x_i:I\longrightarrow\underline{\mathrm{Fl}}(\mathrm{Cat})$ commute aux $\underrightarrow{\lim}$ en question.

3\textsuperscript{o}) Les $x_i$ sont dans $\underline{\mathrm{Fl}}(\mathrm{Cat})_\kappa\overset{\mathrm{def}}{=}\underline{\mathrm{Fl}}_\kappa(\mathrm{Cat})$.

% @ lower third, lines 16–20
Notons que dans le cas d'un ``\unsure{paratopos|prétopos} \unsure{paradisien|paradislien|paradisiaque}'' tel que $\mathrm{Cat}$ ou $\underline{\mathrm{Fl}}\,\mathrm{Cat}$, il y a une façon canonique, $\pi$ étant donné, de représenter un objet $f$ comme $\underrightarrow{\lim}$ d'un syst. inductif $L_1$-$\kappa$-adapté.
% @ foot, last two lines
Ici, \deleted{$\unsure{f}$}
\[
f:X\longrightarrow Y
\]
étant donné, on considère l'ens $I$ des
% archive sheet number: 352
\scanpage{354}
% @ head of sheet
\mspage{78}
% @ top, lines 1–3
couples $(X_i\subset X,Y_i\subset Y)$ de sous-catégories de $X$ et de $Y$, telles que l'on ait $X_i,Y_i\in\mathrm{Ob}\,\mathrm{Cat}_{\unsure{\kappa|\pi}}$ (i.e. $\operatorname{card}(X_i),\operatorname{card}(Y_i)\leq\kappa$), et que $f(\deleted{\unsure{X}}X_i)\subset Y_i$.
% @ upper middle, lines 4–6
Il est immédiat que $\deleted{\unsure{\pi}}I$ est \unsure{filtrante|filtrant}, pour la relation d'ordre évidente
\[
(X_i,Y_i)\subset(X_j,Y_j)\quad\Longleftarrow\quad X_i\subset Y_i\text{ et }X_j\subset Y_j,
\]
% @ centre, lines 7–12
et $=$ \unsure{grde|grande} devant $\kappa$, que $I$ est stable par $\underrightarrow{\lim}$ filtrantes de cardinal $\leq\kappa$, et que le foncteur
\[
i\longmapsto(f_i:X_i\unsure{\mathbin{-}|\to}Y_i):I\longrightarrow\underline{\mathrm{Fl}}\,\mathrm{Cat}
\]
commute aux dites $\underrightarrow{\lim}$, enfin les $f_i$ sont dans $\underline{\mathrm{Fl}}(\mathrm{Cat})_\kappa$ par construction.
% @ lower middle, lines 13–15
\deleted{Il est \unsure{évidt|évident|clair}} Ainsi on a un système inductif $L_1$-$\kappa$-adapté, et \inserted{il est clair que} sa limite inductive est can. is. à $f$. \unsure{Ceci vu}

% @ lower third, corollary and cancellations
\inserted{\uuline{Corollaire} 5.3.2} \deleted{\unsure{vu}, le fait que} La condition : que $W$ soit $\pi$-accessible signifie donc que pour \unsure{toute} $f\inserted{{:}\,X\to Y}\in W$, \deleted{\unsure{si on l'ens} \unsure{ordonné}} et pour \unsure{tte} sous-catégorie $\unsure{X_i|X'}$ de $X$ et $Y_i$ de $Y$, de cardinal $\leq\pi$ (avec $f(X_i)\subset Y_i$, mais en fait c'est inutile, quitte à agrandir $\unsure{Y_i|X_i,Y_i}$), il existe \unsure{alors|deux} sous-catégories
% archive sheet number: 353
\scanpage{355}
% @ head of sheet
\mspage{79}
% @ top, lines 1–3
$X_j\supset X_i$ et $Y_j\supset Y_i$ de cardinal $\leq\unsure{\kappa|\pi}$, telles que $f$ induise $f_j:X_j\unsure{\mathbin{-}|\to}Y_j$ et que l'on ait $f_j\in W$.

% @ upper middle, lines 4–11
La notion d'\unsure{accessibilité} de $W$ me semble aussi bien comprise que celle des limites, je n'y vois plus aucun mystère. D'autre part, il me semble que tous les localiseurs fondamentaux utilisés à ce jour y satisfont, (tout comme à l'axiome $\mathrm{Loc}\,8$). \deleted{Pour} Ici, je vais me borner à établir :

% @ lower middle, theorem heading and statement
\uline{\deleted{Proposition}} \inserted{Théorème} 5.4. Le localiseur \unsure{fond} $\mathbb W_\infty$ (équivalences faibles ``\unsure{habituelles|traditionnelles}) satisfait à l'axiome $\mathrm{Loc}(9)$ (en fait, il satisfait à \uline{tous} les axiomes \unsure{vus|vu} jusqu'\unsure{à} présent, savoir $\mathrm{Loc}(1\mathrm{ter},2,3\mathrm{bis},4,5,6,7,8,9)$).
% @ lower left, sloping note beside Theorem 5.4
% slanted
\margin{Mérite bien le nom de \uline{théorème}, surtout si on y inclut \deleted{\unsure{toute}} la liste des axiomes $\mathrm{Loc}$ mentionnés ci-contre ! Ce théorème n'a rien \deleted{\unsure{de}} d'\unsure{évident}, à partir de la définition \unsure{charabialgique|charabialogique|charabiologique} de $\mathbb W_\infty$ -- c'est un résultat (``bien connu'') qui me paraît profond.}
% @ lower third, small lines below the axiom list
Plus précisément, $\mathbb W_\infty$ est $\unsure{\aleph|\mathbb N}_0$-accessible (où $\aleph_0$ est le premier cardinal infini).
% @ foot, last three lines
Je ne reviens pas \uline{ici} sur la vérification des autres axiomes, et me borne à prouver $\mathrm{Loc}(9)$, l'accessibilité, en utilisant ce qui peut m'être utile. Je me
% archive sheet number: 354
\scanpage{356}
% @ head of sheet
\mspage{80}
% @ top, lines 1–4
borne à indiquer la démonstration (assez technique) dans ses grandes lignes. Pour \inserted{le moment, $W$ est à nouveau un loc fond. \unsure{quelcon.|quelconque}, sauf qu'il satisfait à $\mathrm{Loc}(8)$}

% @ upper third, Lemma 5.5 a
\uline{Lemme} 5.5. a) Pour que le \deleted{\unsure{condition}} \inserted{critère} 5.3.1. b) de $\unsure{\kappa|\pi}$-accessibilité de $W$ soit satisfait, il f. et il s. qu'il le soit lorsque les $f_i$ sont $W_{\unsure{\omega|W|w}}$-parfaits.
% @ upper left, beside the lemma
% slanted
\margin{On suppose que $W$ satisfasse à $\mathrm{Loc}(8)$}
% @ middle, Lemma 5.5 b
b) Supposons que $\mathrm{Parf}_{W_\omega}\cap W\subset W^{\unsure{u|n}}$, i.e. que $W$ satisfait à $\mathrm{Loc}(7\,D_{W_\omega})$ - (la plus faible des variantes de l'axiome des $W$-fibrations, à l'exception seulement de $\mathrm{Loc}(7\,\unsure{F|E})$). \deleted{Alors} \inserted{Supposons de plus (pour être tranquille) que $W$ soit} $W$ est $\pi$-accessible \unsure{sss|ssi} l'ens. $\mathrm{Asph}_W\subset\mathrm{Ob}\,\mathrm{Cat}$ des objets $W$-asphériques de $\mathrm{Cat}$ \deleted{\unsure{est}}soit $\deleted{\unsure{W}}\pi$-accessible. \deleted{\unsure{Il est}} Ou \unsure{encor|encore}, que \deleted{\unsure{pour}} la partie $\mathrm{Asph}_W\subset\mathrm{Ob}\,\mathrm{Cat}_{0\text{-conn}}$ soit $\kappa$-accessible, où $\mathrm{Cat}_{0\text{-conn}}$ désigne la sous-catégorie pleine de $\mathrm{Cat}$ formée des objets $0$-connexes.

% @ foot, proof outline, last six lines
Pour l'essentiel, la démonstration est la même que celle de 5.1, en un peu plus technique. L'outil essentiel est toujours la factorisation fonctorielle $f=pi$ \deleted{avec} d'une flèche de $\mathrm{Cat}$, avec $i\in W_\omega$, $p\in\mathrm{Parf}_{W_\omega}$, et la propriété de
% archive sheet number: 355
\scanpage{357}
% @ head of sheet
\mspage{81}
% @ top, lines 1–2
commutation aux $\underrightarrow{\lim}$ filtrantes de la dite factorisation.

% @ upper third, lines 3–8
Revenons maintenant au cas de $W=\mathbb W_\infty$. On doit prouver que $\mathrm{Asph}_{\mathbb W_\infty}\subset\mathrm{Ob}\,\mathrm{Cat}_{0\text{-conn}}$ est $\unsure{\pi|\kappa}_0$-accessible. Donc \deleted{le critère}, prouver le critère correspondant à 5.3.1 b), en se plaçant maintenant dans $\mathrm{Cat}_{0\text{-}}$ \unsure{ald} $\mathrm{Fl}\,\mathrm{Cat}$. On a donc
% @ centre, limit and conditions
\[
X=\underrightarrow{\lim}_{i\in I}X_i
\]
\unsure{repr} $L_1$-$\pi_0$-admissible de $X$ $\mathbb W_\infty$-asph, à prouver que pour tt $i$ dans $I$, $\exists j$ dans $I$ et $i\unsure{\mathbin{-}|\to}j$, avec $X_j$ $\mathbb W_\infty$-asph. \deleted{Pour ceci, on choisit} \inserted{$\big[$J'admets ici les lemmes. \unsure{AQT|AGT}}

% @ lower middle, Lemma 5.6
\uline{Lemme} 5.6. Soit $\mathbb W_0=\{f\in\deleted{\unsure{f}}\mathrm{Fl}\,\mathrm{Cat}\mid \pi_0(f)\ \text{iso}\}$. Alors $\mathbb W_0$ est \deleted{\unsure{Loc}} $\pi_0$-accessible \inserted{(et \unsure{$f_1$-stable|$f_0$-stable|$f_\infty$-stable})}. A fortiori, l'ensemble des objets $0$-connexes de $\mathrm{Cat}$ est $\pi_0$-accessible.

% @ foot, last five lines
Ceci nous ramène, quitte à remplacer $I$ par une cat. cofinale, à supposer que les $X_i$ sont $0$-connexes. (Mais j'ai oublié que cette réduction est déjà faite ! Donc il faut voir 5.6 comme ingrédient
% archive sheet number: 356
\scanpage{358}
% @ head of sheet
\mspage{82}
% @ head, continuation of bracket
de la démonstration de 5.5 b) $\big]\!\big]$

% @ upper third, lines 2–6
Soit $x$ un objet (point) de $X$, $x_i$ son image dans $X_i$. Considérons (via \deleted{\unsure{la}}\unsure{ma|la} construction \unsure{standart|standard} des catégories des chemins) \deleted{\unsure{qu'on} peut itérer \unsure{fonct.}} les $\Omega^1(X,x)$, $\Omega^1(X_i,x_i)$. On trouve
\[
\Omega^1(X,x)\simeq\underrightarrow{\lim}_{I}\Omega^1(X_i,x_i),
\]
% @ centre, lines 7–11
et on vérifie que c'est là encore un système inductif $L_1$-$\unsure{\pi|\kappa}_0$-adapté. En itérant, on trouve pour $n\in\mathbb N^*$
\[
\Omega^n(X,x)\simeq\underrightarrow{\lim}_{I}\Omega^n(X_i,x_i),
\]
limite inductive $L_1$-$\pi_0$-adaptée. Utilisant 5.6, on en conclut que pour
% @ lower middle, lines 13–16
tout $n$, l'e\deleted{\unsure{ns}} sous-ens (supposons $I$ un ens. ordonné, que \deleted{\unsure{di}} \unsure{diable}, ça suffit !)
\[
I_n=\{i\in I\mid \Omega^n(X_i,x_i)\text{ est }0\text{-connexe}\}
\]
est cofinal dans $I$, de plus il est stable par $\underrightarrow{\lim}$ filtrantes de cardinal $\leq\pi_0$.
% @ foot, intersection and reference
Donc
\[
I_\infty\overset{\mathrm{def}}{=}\bigcap I_n
\]
est aussi cofinal dans $I$ (cf $\overline{\underline{\mathrm{XVIII}}}$ 13.10.2, \unsure{p.|n.|v.} 302).
% archive sheet number: 357
\scanpage{359}
% @ head of sheet
\mspage{83}
% @ top, lines 1–5
Or $i\in I_\infty$ signifie que tous les $\Omega^n(X_i,x_i)$ sont $0$-connexes, ce qui signifie aussi que $X_i$ est $\mathbb W_\infty$-asphérique (puisqu'il est déjà $0$-connexe) - (cf $\overline{\underline{\mathrm{XII}}}$ \unsure{cor.} p. 190). Cela achève la démonstration du th. 5.4.

% @ upper third, section heading
\deleted{\unsure{5}}6. \uline{\deleted{\unsure{L'axiome de descente} $\mathrm{Loc}(\unsure{10})$}} \inserted{\uline{La condition} $\mathrm{Loc}(5,D_{W_{\unsure{\omega|W|w}}})$ \uline{implique} $W\subset\mathbb W_\infty$.}

% @ centre, introduction and theorem
Je vais redonner ici, pour mémoire, la démonstration du th. 1 de $\overline{\underline{\mathrm{XII}}}$ :

\uline{Théorème} 6.1. Soit $W$ un localiseur fondamental satisfaisant à $D_{W_\omega}$ (où $W_\omega$ est le plus petit de tous les localiseurs fondamentaux). Alors on a $W=\mathrm{Fl}(\mathrm{Cat})$, \uline{ou} $W\subset\mathbb W_\infty$ (suivant que $W$ ne satisfait pas à l'axiome de connexité $\mathrm{Loc}(5)$, ou qu'il y satisfait).

% @ lower third, references following the theorem
Je vais faire la démonstration en supposant que $W$ satisfait à $\mathrm{Loc}(5)$. Pour la démonstration du fait que dans le cas contraire, on a $W=\mathrm{Fl}(\mathrm{Cat})$, je renvoie à $\overline{\underline{\mathrm{XII}}}$ th. 2 p. 20. (Voir aussi le résultat \inserted{plus général et plus faible,} très voisin, \deleted{\unsure{p. 12}} \inserted{rappelé \unsure{en}} prop 2.7. page 12.)

% @ foot, Lemma 6.2
\uline{Lemme} 6.2. Soient $W,W'$ deux loc fond., \inserted{$W$} satisfaisant à $\mathrm{Loc}(7\,D_{W_\omega})$. Pour qu'on ait $W\subset W'$, il suffit qu'on ait $\mathrm{Asph}_W\subset\mathrm{Asph}_{\unsure{W_\infty|W'}}$
% archive sheet number: 358
\scanpage{360}
% @ head of sheet
\mspage{84}
% @ top, lines 1–6
Il suffit bien sûr de prouver ``il suffit'', donc supposons que tout objet $W$-asphérique est $W'$-asphérique, de prouver que $W\subset W'$. Soit donc $f:X\unsure{\mathbin{-}|\to}Y$ dans $W$, prouvons $f\in W'$. Pour ceci, j'u\deleted{\unsure{f}}tilise la factorisation canonique $f=pi$
% @ upper middle, factorization
\diagram{\begin{tikzcd}[ampersand replacement=\&]
X \arrow[r,"i\in W_{\unsure{\omega|W|w}}"] \& \overline X \arrow[r,"p\in\mathrm{Parf}_{W_\omega}"] \& Y{.}
\end{tikzcd}}
% @ middle, lines 8–16
On a donc aussi $i\in W$, et comme $f=pi\in W$, on a $p\in W$ donc $p\in\mathrm{Parf}_{W_\omega}\cap W$. Or $D_{W_\omega}$ signifie que cet ensemble est $\subset W^{\unsure{u|n}}$, donc on a $p\in W^u$, \deleted{\unsure{et}} en particulier les fibres de $p$ sont $W$-asphériques. Elles sont donc $W'$-asphériques, et comme $p$ est $W_\omega$-parfait et a fortiori $W'$-parfait, il s'ensuit que $p\in W'$. On a aussi $i\in W_\omega\subset W'$, donc $i\in W'$, donc aussi $f=pi\in W'$, \unsure{qed.|qcd.}

% @ lower third, last seven lines
Revenons donc à la dém. de 6.1. En vertu du lemme, il suffit de prouver que si $X$ dans $\mathrm{Cat}$ est $W$-asphérique, il est $\mathbb W_\infty$-asphérique. Mais par hyp. $\mathrm{Loc}(5)$ on sait déjà que $X$ est $0$-connexe, et choisissant un pt $x$ dans $X$, qu'il soit $\mathbb W_\infty$-asphérique s'exprime par la condition
% archive sheet number: 359
\scanpage{361}
% @ head of sheet
\mspage{85}
% @ top, lines 1–7
que les $\Omega^n(X,x)$ (pour $n\geq1$) sont $0$-connexes. (Cf. $\overline{\underline{\mathrm{XII}}}$ cor. p. 190 déjà cité.) Pour ceci, à cause de $\mathrm{Loc}(5)$ il suffit de prouver qu'ils sont $W$-asphériques, et par induction il suffit de prouver que $\Omega^1(X,x)$ l'est. Or par définition de $\Omega^1(X,x)$, c'est là la fibre en $x$ d'une flèche dans $\mathrm{Cat}$
% @ upper middle, path-category arrow
\diagram{\begin{tikzcd}[ampersand replacement=\&]
E=\underline{\mathrm{Ch}}(X,x,-) \arrow[r,"p\in\mathrm{Parf}_{W_{\unsure{\omega|W|w}}}"] \& X{,}
\end{tikzcd}}
% @ middle, lines 9–14
où la source $E$ est $W_\omega$-asphérique, a fortiori $W$-asphérique. Comme $X$ l'est également, on a donc $p\in W$, donc $p\in\mathrm{Parf}_{W_\omega}\cap W\subset W^{\unsure{u|n}}$, donc $p$ est à fibres $W$-asphériques, et en particulier $\Omega^1(X,x)$ est $W$-asphérique, qed.

% @ lower third, Corollary 6.3
\uline{Corollaire} 6.3. (Caractérisation de $\mathbb W_\infty$). $\mathbb W_\infty$ est le plus grand localiseur fondamental satisfaisant à $\mathrm{Loc}(5,7\,D_{W_\omega})$ - donc aussi le plus grand loc. fond. satisfaisant à $\mathrm{Loc}(5,7)$.

% @ foot, final three lines
Je pense que nous trouverons plus bas une caractérisation ``interne'' de $\mathbb W_\infty$, via un ``axiome de descente'' (\deleted{\unsure{b}} \unsure{très} fort) $\mathrm{Loc}(10)$,
% archive sheet number: 360
\scanpage{362}
% @ head of sheet
\mspage{86}
% @ top, conclusion of hopes question
comme l'unique localiseur fond. satisfaisant $\mathrm{Loc}(5,7D_{W_{\unsure{\omega|W|w}}},10)$, (ou encor $\mathrm{Loc}(5,7,10)$).

% @ upper third, heading 6.4
6.4. \uline{Commentaires sur les axiomes des} $\unsure{W|\mathbb W}$-\uline{fibrations du type} $\mathrm{Loc}(7\ldots)$, \uline{et notamment sur} $\mathrm{Loc}(7D_{W_\omega})$. Relation \inserted{conjecturale} avec les ``\uline{types d'homotopie \unsure{simples|simple}}'' de Whitehead.

% @ middle, a), through the comparison with the theorem
a) Il me semble que \uline{tous} les localiseurs fondamentaux avec lesquels on a \unsure{travaillé} tacitement jusqu'à présent (les $\mathbb W_n$ $\bigl(n\in\mathbb N^{\unsure{\alpha|\infty}}\bigr)$, équivalences homologiques et cohomologiques, équivalences rationnelles $\ldots$) correspondent à des localiseurs $W\supset\mathbb W_\infty$. Par le th. 6.1, il s'ensuit donc qu'\uline{aucun} d'eux, à part $\mathbb W_\infty$ lui-même, ne satisfait à \inserted{la} plus faible version \inserted{$D_{W_\omega}$} de toutes de l'axiome $\mathrm{Loc}(7\ldots)$, \deleted{\unsure{à|sauf}} \unsure{à} la seule \deleted{\unsure{seule}} exception de $\mathrm{Loc}(7F)$. Ainsi, les localiseurs $W$ qui y satisfont \deleted{et qui} \inserted{sont} \deleted{ne contiennent \unsure{donc} pas} \inserted{à l'exception de $D_{W_\omega}$ lui-même,} d'un type \deleted{tout} nouveau -- des localiseurs qui sont $\subsetneq\mathbb W_\infty$.

% @ foot, b), continues
b) Ceci soulève, bien sûr, la question de l'existence de loc. fond. $\subset\mathbb W_\infty$ et
% archive sheet number: 361
